%-----------------------------------LICENSE------------------------------------%
%   This file is part of Mathematics-and-Physics.                              %
%                                                                              %
%   Mathematics-and-Physics is free software: you can redistribute it and/or   %
%   modify it under the terms of the GNU General Public License as             %
%   published by the Free Software Foundation, either version 3 of the         %
%   License, or (at your option) any later version.                            %
%                                                                              %
%   Mathematics-and-Physics is distributed in the hope that it will be useful, %
%   but WITHOUT ANY WARRANTY; without even the implied warranty of             %
%   MERCHANTABILITY or FITNESS FOR A PARTICULAR PURPOSE.  See the              %
%   GNU General Public License for more details.                               %
%                                                                              %
%   You should have received a copy of the GNU General Public License along    %
%   with Mathematics-and-Physics.  If not, see <https://www.gnu.org/licenses/>.%
%------------------------------------------------------------------------------%
\documentclass{article}
\usepackage{graphicx}
\usepackage{amsmath, amsthm, amssymb}
\usepackage{xcolor}
\theoremstyle{normal}
\newtheorem{problem}{Problem}
\newtheorem{theorem}{Theorem}
\title{The Arzel\'{a}-Ascoli Theorem}
\author{Ryan Maguire}
\date{\today}
\setlength{\parindent}{0em}
\setlength{\parskip}{0em}
\newif\ifsolution
\solutiontrue
\graphicspath{{../../../images/}}

\begin{document}
    \maketitle
    \textbf{Continuity} may be written as follows.
    Given $f:A\rightarrow\mathbb{R}$, with $A\subseteq\mathbb{R}$,
    $f$ is continuous provided
    \begin{equation}
        \forall_{x_{0}\in{A}}
        \forall_{\varepsilon>0}
        \exists_{\delta>0}
        \forall_{x\in{A}}
        \Big(
            |x-x_{0}|<\delta
            \Rightarrow
            |f(x)-f(x_{0})|<\varepsilon
        \Big),
    \end{equation}
    which reads
    \textit{for all $x_{0}\in{A}$ and for all $\varepsilon>0$}
    \textit{there exists a $\delta>0$ such that for all $x\in{A}$},
    \textit{if $|x-x_{0}|<\delta$, then $|f(x)-f(x_{0})|<\varepsilon$}.
    Writing the definition in the former way provides few benefits since it is
    very hard to read, but since our eyes travel from left-to-right across
    the page, this expression helps us see which variables are allowed to
    depend on which other variables. In particular, $x_{0}\in{A}$ is chosen
    first, and it may be anything in $A$. Then $\varepsilon>0$ is chosen, and
    it may be any positive real number. The statement \textit{there exists}
    $\delta>0$ occurs after $x_{0}$ and $\varepsilon$ have been chosen, and so
    it is allowed to depend on these two variables. We may write
    \begin{equation}
        \delta
        =
        \delta(x_{0},\,\varepsilon).
    \end{equation}
    \textbf{Uniform continuity} is defined by slightly modifying the
    definition of continuity. We have
    \begin{equation}
        \forall_{\varepsilon>0}
        \exists_{\delta>0}
        \forall_{x_{0}\in{A}}
        \forall_{x\in{A}}
        \Big(
            |x-x_{0}|<\delta
            \Rightarrow
            |f(x)-f(x_{0})|<\varepsilon
        \Big).
    \end{equation}
    All we have done is move $\forall_{x_{0}\in{A}}$ to occur \textit{after}
    the statement for $\delta$. Hence $\delta$ is \textit{not} allowed to
    depend on $x_{0}$, so we should write
    \begin{equation}
        \delta
        =
        \delta(\varepsilon).
    \end{equation}
    A small difference, but an important one.
    We proved in class that if $A$ is compact, and if $f:A\rightarrow\mathbb{R}$
    is continuous, then $f$ is uniformly continuous.
    \par\hfill\par
    Similar rearrangements of quantifiers occur with other notions.
    If $C(\mathbb{R},\,\mathbb{R})$ denotes the set of all continuous
    functions $f:\mathbb{R}\rightarrow\mathbb{R}$, then given a sequence
    of functions $F:\mathbb{N}\rightarrow{C}(\mathbb{R},\,\mathbb{R})$,
    we say that $F_{n}$ \textbf{converges point-wise} to a function $f$
    provided
    \begin{equation}
        \forall_{x_{0}\in\mathbb{R}}
        \forall_{\varepsilon>0}
        \exists_{N\in\mathbb{N}}
        \forall_{n\in\mathbb{N}}
        \Big(
            n>N
            \Rightarrow
            |f(x_{0})-F_{n}(x_{0})|<\varepsilon
        \Big).
    \end{equation}
    Note that $x_{0}$ occurs \textit{before} $N$, and hence $N$ is allowed to
    depend both on $\varepsilon$ and $x_{0}$, so we write
    \begin{equation}
        N
        =
        N(x_{0},\,\varepsilon).
    \end{equation}
    \textbf{Uniform convergence} occurs when $N$ may be chosen independent of
    $x_{0}$. That is:
    \begin{equation}
        \forall_{\varepsilon>0}
        \exists_{N\in\mathbb{N}}
        \forall_{x_{0}\in\mathbb{R}}
        \forall_{n\in\mathbb{N}}
        \Big(
            n>N
            \Rightarrow
            |f(x_{0})-F_{n}(x_{0})|<\varepsilon
        \Big).
    \end{equation}
    We then write
    \begin{equation}
        N
        =
        N(\varepsilon).
    \end{equation}
    We proved that if $F_{n}$ is continuous for each $n$, and
    $F_{n}\rightarrow{f}$ uniformly, then $f$ is continuous as well.
    \begin{problem}
        \par\hfill\par
        Prove that if $A$ is compact, if $F_{n}:A\rightarrow\mathbb{R}$ is
        continuous for all $n\in\mathbb{N}$, and if $F_{n}\rightarrow{f}$
        uniformly, then the sequence $F$ is uniformly bounded. That is, there
        is an $M\in\mathbb{R}$ such that for all $n\in\mathbb{N}$ and for all
        $x\in{A}$ we have $|F_{n}(x)|<M$.
    \end{problem}
    \ifsolution
        \color{blue}
        \begin{proof}[Solution]
            Since $F_{n}\rightarrow{f}$ uniformly, and since each $F_{n}$ is
            continuous, we know that $f$ is continuous. But $A$ is compact
            and hence by the extreme value theorem, $f$ is bounded. Let
            $M_{f}>0$ be a bound and let $\varepsilon=1$. Since $\varepsilon>0$
            and $F_{n}\rightarrow{f}$ uniformly, there is an $N\in\mathbb{N}$
            such that for all $x_{0}\in{A}$ and for all $n>N$ we have
            $|F_{n}(x_{0})-f(x_{0})|<\varepsilon$. Since each function
            $F_{0}$, $F_{1}$, $\dots$, $F_{N}$ is continuous and $A$ is compact,
            each of these is bounded. Let $M_{0}$, $M_{1}$, $\dots$, $M_{N}$
            be bounds for $F_{0}$, $F_{1}$, $\dots$, $F_{N}$, respectively.
            Define $M$ to be:
            \begin{equation}
                M
                =
                \textrm{max}\left(
                    M_{0},\,M_{1},\,\cdots,\,M_{N},\,M_{f}
                \right)+1.
            \end{equation}
            Then for all $x_{0}\in{A}$ and for all $n>N$ we have:
            \begin{equation}
                \begin{array}{rcl}
                    \displaystyle
                    |F_{n}(x_{0})|
                    \!\!
                    &=&
                    \!\!
                    \displaystyle
                    |F_{n}(x_{0})-f(x_{0})+f(x_{0})|\\[1em]
                    \!\!
                    &\leq&
                    \!\!
                    \displaystyle
                    |F_{n}(x_{0})-f(x_{0})|+|f(x_{0})|\\[1em]
                    \!\!
                    &<&
                    \!\!
                    \displaystyle
                    \varepsilon+M_{f}\\[1em]
                    \!\!
                    &=&
                    \!\!
                    \displaystyle
                    1+M_{f}\\[1em]
                    \!\!
                    &<&
                    \!\!
                    \displaystyle
                    M
                \end{array}
            \end{equation}
            And $|F_{n}(x_{0})|<M$ by definition for each $n\leq{N}$, so
            $F$ is uniformly bounded.
        \end{proof}
        \color{black}
    \fi
    \begin{problem}
        Prove that if $A$ is compact, if $F_{n}:A\rightarrow\mathbb{R}$ is
        continuous for all $n\in\mathbb{N}$, and if
        $F_{n}\rightarrow{f}$ uniformly, then $F$ is
        \textbf{uniformly equicontinuous}, meaning for all $\varepsilon>0$
        there is a $\delta>0$ such that for all $n\in\mathbb{N}$, for all
        $x_{0}\in{A}$, and for all $x\in{A}$, if
        $|x-x_{0}|<\delta$, then
        $|F_{n}(x)-F_{n}(x_{0})|<\varepsilon$.
        That is
        \begin{equation}
            \forall_{\varepsilon>0}
            \exists_{\delta>0}
            \forall_{n\in\mathbb{N}}
            \forall_{x_{0}\in{A}}
            \forall_{x\in{A}}
            \Big(
                |x-x_{0}|<\delta
                \Rightarrow
                |F_{n}(x)-F_{n}(x_{0})|<\varepsilon
            \Big).
        \end{equation}
        Note that this means that $\delta$ may be chosen independent of
        $x_{0}$ and independent of $n$.
    \end{problem}
    \ifsolution
        \color{blue}
        \begin{proof}[Solution]
            For let $\varepsilon>0$. Since $A$ is compact, and since
            each $F_{n}$ is continuous, we may conclude that each $F_{n}$ is
            uniformly continuous. Similarly, since $F_{n}\rightarrow{f}$
            uniformly, we may conclude that $f$ is continuous, and since
            $A$ is compact, $f$ must be uniformly continuous as well.
            Hence for each $n\in\mathbb{N}$ there is a $\delta_{n}>0$ such that
            for all $x,x_{0}\in{A}$, if $|x-x_{0}|<\delta$, then
            $|F_{n}(x)-F_{n}(x_{0})|<\varepsilon/3$. Similarly, there is a
            $\delta_{f}$ for $f$. Since $F_{n}\rightarrow{f}$ uniformly, there
            is an $N\in\mathbb{N}$ such that for all $x\in{A}$ and for all
            $n>N$ we have $|f(x)-F_{n}(x)|<\varepsilon/3$. Define $\delta$
            by
            \begin{equation}
                \delta
                =
                \min\left(
                    \delta_{0},\,
                    \delta_{1},\,
                    \dots,\,
                    \delta_{N},\,
                    \delta_{f}
                \right).
            \end{equation}
            If $n\leq{N}$, then by definition of $\delta$ we have
            $|x-x_{0}|<\delta$ implies
            $|F_{n}(x)-F_{n}(x_{0})|<\varepsilon$. If $n>N$, then for all
            $x,x_{0}\in{A}$ with $|x-x_{0}|<\delta$, we have:
            \begin{equation}
                \begin{array}{rcl}
                    &&
                    \!\!
                    \displaystyle
                    |F_{n}(x)-F_{n}(x_{0})|\\[1em]
                    \!\!
                    &=&
                    \!\!
                    \displaystyle
                    |F_{n}(x)-f(x)+f(x)-f(x_{0})+f(x_{0})-F_{n}(x_{0})|\\[1em]
                    \!\!
                    &\leq&
                    \!\!
                    \displaystyle
                    |F_{n}(x)-f(x)|
                    +
                    |f(x)-f(x_{0})|
                    +
                    |f(x_{0})-F_{n}(x_{0})|\\[1em]
                    \!\!
                    &<&
                    \!\!
                    \displaystyle
                    \frac{\varepsilon}{3}
                    +
                    \frac{\varepsilon}{3}
                    +
                    \frac{\varepsilon}{3}\\[1em]
                    \!\!
                    &=&
                    \!\!
                    \varepsilon,
                \end{array}
            \end{equation}
            and hence $F$ is uniformly equicontinuous.
        \end{proof}
        \color{black}
    \fi
    The other direction works as well.
    \begin{theorem}[Arzel\'{a}-Ascoli Theorem]
        If $F:\mathbb{N}\rightarrow{C}(A,\,\mathbb{R})$ is a uniformly
        equicontinuous sequence of functions on a compact set $A$, and if
        $F_{n}\rightarrow{f}$, then the convergence is uniform.
    \end{theorem}
    \begin{proof}
        To show that $F_{n}\rightarrow{f}$ uniformly, we show that $F$ defines
        a \textit{uniformly Cauchy} sequence. That is, $F_{n}(x)$ is Cauchy
        for each $x\in{A}$, and the $N\in\mathbb{N}$ such that
        $n,m>N$ implies $|F_{n}(x)-F_{m}(x)|<\varepsilon$ may be chosen
        independent of $x$.
        \par\hfill\par
        Let $\varepsilon>0$.
        Since $F$ is uniformly equicontinuous, there is a $\delta>0$ such
        that for all $x,x_{0}\in{A}$, and for all $n\in\mathbb{N}$, if
        $|x-x_{0}|<\delta$, then $|F_{n}(x)-F_{n}(x_{0})|<\varepsilon/3$.
        Since $F_{n}(x)\rightarrow{f}(x)$ (point-wise), for each
        $x\in{A}$ there is an $N_{x}\in\mathbb{N}$ such that
        $n,m>N_{x}$ implies $|F_{m}(x)-F_{n}(x)|<\varepsilon/3$.
        That is, point-wise convergence implies point-wise Cauchy sequences.
        Now consider the open cover $\mathcal{O}$ given by
        \begin{equation}
            \mathcal{O}
            =
            \left\{\,
                \left(x-\frac{\delta}{4},\,x+\frac{\delta}{4}\right)
                \;\Big|\;
                x\in{A}\,
            \right\}.
        \end{equation}
        $\mathcal{O}$ is an open cover of $A$ since
        $x\in(x-\delta/4,\,x+\delta/4)$ for each $x\in{A}$.
        Since $A$ is compact, there is a finite subcover $\Delta$ which we
        may write as
        \begin{equation}
            \Delta
            =
            \left\{\,
                \left(
                    x_{0}-\frac{\delta}{4},\,x_{0}+\frac{\delta}{4}
                \right),\,
                \dots,\,
                \left(
                    x_{k}-\frac{\delta}{4},\,x_{k}+\frac{\delta}{4}
                \right)
            \right\}.
        \end{equation}
        For each $x_{0}$, $x_{1}$, $\dots$, $x_{k}$ there is an integer
        $N_{0}$, $N_{1}$, $\dots$, $N_{k}$ attached to it from before.
        Define $N=\textrm{max}(N_{0},\,\dots,\,N_{k})$, and suppose
        $n,m>N$. If $x\in{A}$, then since $\Delta$ is an open cover
        of $A$, there is some $x_{M}\in{A}$ such that
        $x\in(x_{M}-\delta/4,\,x_{M}+\delta/4)$. But then:
        \begin{equation}
            \begin{array}{rcl}
                &&
                \!\!
                \displaystyle
                |F_{n}(x)-F_{m}(x)|\\[1em]
                \!\!
                &=&
                \!\!
                \displaystyle
                |
                    F_{n}(x)-F_{n}(x_{M})
                    +
                    F_{n}(x_{M})-F_{m}(x_{M})
                    +
                    F_{m}(x_{M})-F_{m}(x)
                |\\[1em]
                \!\!
                &\leq&
                \!\!
                \displaystyle
                |F_{n}(x)-F_{n}(x_{M})|
                +
                |F_{n}(x_{M})-F_{m}(x_{M})|
                +
                |F_{m}(x_{M})-F_{m}(x)|.
            \end{array}
        \end{equation}
        Now $|F_{n}(x)-F_{n}(x_{M})|<\varepsilon/3$ since
        $|x-x_{M}|<\delta$, and similarly
        $|F_{m}(x)-F_{m}(x_{M})|<\varepsilon/3$.
        Finally, $|F_{n}(x_{M})-F_{m}(x_{M})|<\varepsilon/3$
        since $n,m>N$ ($F_{n}(x_{M})$ is a Cauchy sequence).
        Thus:
        \begin{equation}
            \begin{array}{rcl}
                &&
                \!\!
                \displaystyle
                |F_{n}(x)-F_{m}(x)|\\[1em]
                \!\!
                &\leq&
                \!\!
                \displaystyle
                |F_{n}(x)-F_{n}(x_{M})|
                +
                |F_{n}(x_{M})-F_{m}(x_{M})|
                +
                |F_{m}(x_{M})-F_{m}(x)|\\[1em]
                \!\!
                &<&
                \!\!
                \displaystyle
                \frac{\varepsilon}{3}
                +
                \frac{\varepsilon}{3}
                +
                \frac{\varepsilon}{3}\\[1em]
                \!\!
                &=&
                \!\!
                \displaystyle
                \varepsilon.
            \end{array}
        \end{equation}
        That is, $F_{n}$ is uniformly Cauchy, and hence
        $F_{n}\rightarrow{f}$ uniformly.
    \end{proof}

\end{document}
